\clearpage
\begingroup
\fancyhead[L]{\small REMARKS}
\renewcommand{\headrulewidth}{0pt}
\par\noindent\hspace*{\parindent}On [VII.], the letter extracts are printed from the originals. The editor
is responsible for the whole of Note [2.]. The enormously extensive
calculations for these tables (see p. 429), together with the completed table itself, form
a packet of neatly arranged sheets. Many of them are in hands other than
\textsc{Gauss}'s, having been written by \textsc{Encke} and \textsc{Westphal}.

\par\noindent\hspace*{\parindent}On [VIII.], among the posthumous papers the material on the perturbations of Pallas by Saturn
contains nothing in \textsc{Gauss}'s hand. It consists solely of the original letters from
\textsc{Nicolai} to \textsc{Gauss}, together with some sheets in \textsc{Nicolai}'s hand.
Equations 1) and 2) in Note [1.] come from these sheets. The remaining
contents of that note were supplied. The relevant letter extracts
are printed in Notes [2.]--[6.]. The results for the perturbations are on
separate sheets that originally accompanied \textsc{Nicolai}'s letters
and are now filed among the boxes containing the Pallas perturbations in the posthumous papers.
The results are printed in a somewhat different order from that of the
original sheets. There they are arranged in Manipel, while the printed
text follows their original natural sequence.

\par\noindent\hspace*{\parindent}The complete calculation of these perturbations survives among \textsc{Nicolai}'s papers
at the Grand Ducal Observatory in Heidelberg, and Mr
\textsc{Valentiner} kindly allowed the undersigned
to inspect it. It bears the inscription »Geendigt den
20\textsuperscript{ten} April 1814.« We may perhaps still hope to see it published
in somewhat greater detail elsewhere. Such a publication would scarcely have been appropriate in \textsc{Gauss}'s Works.
\textsc{Nicolai} also carried out two calculations of the
Saturn perturbations. The first dealt only with the principal equations,
and its results are likewise preserved among the posthumous papers. The corresponding
letters from \textsc{Gauss} to \textsc{Nicolai} appear to have been lost. Mr
\textsc{Valentiner} at any rate informed me that \textsc{Nicolai}'s papers contain only
a series of letters from \textsc{Gauss} from a later period (from 1819 onward),
which Mr \textsc{Valentiner} has published (Briefe von C. F.
\textsc{Gauss} an B. \textsc{Nicolai}, Karlsruhe 1877). The letters from \textsc{Nicolai}
to \textsc{Gauss} appear to have survived in full and have yet to
be published.

\par\noindent\hspace*{\parindent}On [IX.], the formulas and numbers in Note [1.] are found partly in
notebook Bc, partly on loose sheets, and partly in a booklet entitled »Astronomische Rechnungen;
Angefangen im März 1814«, which otherwise deals with other subjects
and is therefore filed among the posthumous papers under »Beobachtungen und Rechnungen«.
The numbers in Notes [2.]--[3.] come from calculations that \textsc{Gauss} carried out on separate sheets
and left unfinished. Note [4.]
is the editor's work. [5.]--[6.] come from the booklet »Astronomische Rechnungen« mentioned above,
with some relations in Note [5.] also coming from a separate sheet. Finally, [7.] comes
from notebook Be, where the passage printed here is followed by a worked example
(for \(M=165^0\)).

\par\noindent\hspace*{\parindent}Notes [4.]--[7.] describe a very interesting
device for expanding the disturbing function. In the end \textsc{Gauss}
set it aside, since it seems to offer little saving
in numerical calculation. By this method
\textsc{Gauss} calculated the printed values for \(M=30^0\) and only
those for \(M=15^0\) and \(45^0\) besides.

\medskip
\par\noindent\hspace*{\parindent}We see that \textsc{Gauss} brought the enormous problem of calculating
the perturbations of Pallas within the limits of observational accuracy almost to completion.
Even today an astronomer would be reluctant to tackle it.
Only a few perturbation equations for Mars are missing.
It was all the more regrettable that this work has remained
unknown to this day. As the last of the printed letter extracts shows,
\textsc{Gauss} himself felt this keenly.

\par\noindent\hspace*{\parindent}It is very striking that \textsc{Gauss} never once mentions this investigation
in his scientific diary (or journal of notes, as he himself calls it), which Mr
F. \textsc{Klein} edited in the »Festschrift zur Feier des 150jährigen Bestehens der K. Ge%
\SourceLeaf{616}sellschaft der Wissenschaften zu Göttingen«, published in
1901.
For the years 1802--1804, the diary gives the following remark. »\textit{Annis
insequentibus 1802, 1803, 1804 occupationes astronomicae maximam otii partem
abstulerunt, calculi imprimis circa planetarum novorum theoriam instituti.
Unde evenit, quod hisce annis catalogus hicce neglectus est} \mbox{.\ .\ .\ .\ .\ .\ .}«
In these years fall \textsc{Gauss}'s studies of the two-body problem
(in preparation for the Theoria motus) and the perturbations of Ceres. It is curious
that \textsc{Gauss}'s engagement with astronomical calculations
always interrupted the keeping of the diary. For in
1810, when \textsc{Gauss} began working on Pallas,
another long gap begins. On 29 February 1812 he then writes »\textit{Catalogum
praecedentem per fata iniqua iterum interruptum initio anni 1812 resumimus}
\mbox{.\ .\ .\ .\ .}« He then mentions rediscovering the second (»zweiten«) proof
of the fundamental theorem of algebra, which he had found in November 1811
and subsequently forgotten in part because he had written nothing
down. Between this entry and the next one of 26 September
1812, which concerns the attraction of an ellipsoid, falls the
discovery of the rationality of the mean motions of Pallas and Jupiter.
The diary is entirely silent about this as well.

\par\noindent\hspace*{\parindent}\textsc{Gauss}'s investigations of the perturbations of Pallas
clearly still provide valuable material today. It would be highly
desirable to reduce the perturbation expressions to the present
epoch.
\par\nobreak\noindent\hfill\textsc{Brendel}.\par
\nobreak\vspace{3pt}\noindent\makebox[\linewidth]{\rule{.24\linewidth}{.4pt}}\par
\clearpage
\endgroup
